%! AP Statistics — Unit 8, Lesson 8.3 — Homework (Answer Key)
\documentclass[10pt]{article}
\usepackage{apstats-article}
\usepackage{apstats-key}

\begin{document}

\pageheader{Unit 8, Lesson 8.3}{Homework --- Answer Key}
\namedateperiod

\vspace{0.1in}
\textbf{Directions:} Show all work for full credit. Write complete sentences for
interpretations and conclusions.

\begin{enumerate}

  \item \textbf{Genetics.}

        \renewcommand{\arraystretch}{1.4}
        \begin{center}
        \begin{tabular}{lcccc}
            \textbf{Phenotype} & 1  & 2  & 3  & 4  \\
            \hline
            \textbf{Observed}  & 89 & 31 & 27 & 13 \\
        \end{tabular}
        \end{center}

        \begin{enumerate}[label=\alph*.]
          \item Expected counts under 9:3:3:1 (total = 16 parts, $n = 160$):
                \ansline{Phenotype 1: $160 \times 9/16 = 90$; Phenotype 2: $160 \times 3/16 = 30$; Phenotype 3: $160 \times 3/16 = 30$; Phenotype 4: $160 \times 1/16 = 10$.}

          \item Chi-square contributions:
                \ansline{$(89 - 90)^2/90 + (31 - 30)^2/30 + (27 - 30)^2/30 + (13 - 10)^2/10 = 1/90 + 1/30 + 9/30 + 9/10 \approx 0.011 + 0.033 + 0.300 + 0.900 = 1.244$}

                \begin{teachernote}
                    Some rounding may yield $\chi^2 \approx 1.24$--$1.25$; accept values in this range.
                    The p-value given in the problem ($\approx 0.81$) corresponds to this statistic.
                \end{teachernote}

          \item $df =$ \ans{$4 - 1 = 3$}

          \item \ansline{Because $p \approx 0.81 > \alpha = 0.05$, we fail to reject $H_0$. We do not have convincing evidence that the phenotype distribution of offspring differs from the 9:3:3:1 ratio predicted by Mendelian genetics.}
        \end{enumerate}

  \item \textbf{Streaming Services.}

        \renewcommand{\arraystretch}{1.4}
        \begin{center}
        \begin{tabular}{lccccc}
            \textbf{Service}  & Netflix & Hulu & Disney+ & Amazon & Other \\
            \hline
            \textbf{Observed} & 78      & 42   & 35      & 28     & 17    \\
        \end{tabular}
        \end{center}

        \begin{enumerate}[label=\alph*.]
          \item \ansline{$H_0$: The distribution of most-used streaming service is equal across all 5 services (each at 20\%). $H_a$: At least one service is used most by a different proportion of college students.}

          \item Expected count for each service: \ans{$200 \times 0.20 = 40$}.

                Chi-square contributions:
                \ansline{$(78-40)^2/40 + (42-40)^2/40 + (35-40)^2/40 + (28-40)^2/40 + (17-40)^2/40 = 1444/40 + 4/40 + 25/40 + 144/40 + 529/40 = 36.1 + 0.1 + 0.625 + 3.6 + 13.225 = 53.65$}

          \item $df =$ \ans{$5 - 1 = 4$}

          \item \ansline{The probability of obtaining a chi-square statistic of 53.65 or larger, given that all five streaming services are equally popular among college students, is less than 0.001. Because $p < 0.001 < \alpha = 0.05$, we reject $H_0$. We have convincing evidence that the distribution of most-used streaming service is not equal across all five services.}
        \end{enumerate}

  \item \textbf{Conceptual Question.}

        \ansline{Failing to reject $H_0$ does not prove $H_0$ is true. It means only that the data are consistent with the null model --- the observed discrepancy could reasonably occur by chance if $H_0$ were true. We may have insufficient evidence to detect a real difference (a Type II error). We can never ``prove'' a null hypothesis through a significance test.}

\end{enumerate}

\end{document}
